\section{为什么线性代数需要复数}
\label{sec:03-Linear-Algebra/08-Complex-Matrices-and-Fourier-Structure/01}

一个二维状态每步按 \(x_{k+1}=Ax_k\) 更新，你想知道它长期会跑到哪里去。标准做法是把初值拆到特征方向上，让各分量按 \(\lambda^k\) 各自缩放。换成这个矩阵试试：

\[
A=\begin{bmatrix}0&-1\\1&0\end{bmatrix}.
\]

它每步把状态逆时针转 \(90^\circ\)。轨迹本身毫无病态——四步一个循环，规规矩矩地在圆上走。可特征多项式是 \(\lambda^2+1\)，实数里没有根，一条不变方向也画不出来。\(v=(a,b)^{\trans}\) 被送到 \((-b,a)^{\trans}\)，两者点积恒为零，永远垂直。

失效的不是这个系统，是我们手上的标量。实数只能表达``沿同一条直线拉长或压短''，而旋转把每条直线都挪了位置，没有一条能用一次实数乘法描述。要让 \(Av=\lambda v\) 在这里仍有解，\(\lambda\) 必须同时携带``转多少''和``缩多少''两条信息。复数恰好是这样的标量：\(z=re^{i\theta}\) 里 \(r\) 管伸缩，\(\theta\) 管转角。

\begin{center}
\begin{tikzpicture}[x=1cm,y=1cm,font=\footnotesize,>=stealth]
  \draw[black!20] (0,0) circle (1.5);
  \draw[->,black!45] (-1.9,0) -- (2.0,0);
  \draw[->,black!45] (0,-1.9) -- (0,2.0);
  \foreach \a in {8,58,118,168,218,268,318} {
    \draw[->,blue!65!black,thick] (0,0) -- ({1.5*cos(\a)},{1.5*sin(\a)});
    \draw[->,red!55!black,dashed] ({1.62*cos(\a)},{1.62*sin(\a)}) arc (\a:\a+50:1.62);
  }
  \node at (0,-2.5) {实平面：没有一条直线留在原处};
  \begin{scope}[shift={(6.2,0)}]
    \draw[black!35] (0,0) circle (1.5);
    \draw[->,black!45] (-1.9,0) -- (2.0,0) node[below right] {\(\operatorname{Re}\)};
    \draw[->,black!45] (0,-1.9) -- (0,2.0) node[left] {\(\operatorname{Im}\)};
    \fill[blue!65!black] ({1.5*cos(50)},{1.5*sin(50)}) circle (2.2pt);
    \node[right] at ({1.5*cos(50)+0.12},{1.5*sin(50)+0.12}) {\(e^{i\theta}\)};
    \fill[blue!65!black] ({1.5*cos(-50)},{1.5*sin(-50)}) circle (2.2pt);
    \node[right] at ({1.5*cos(-50)+0.12},{1.5*sin(-50)-0.12}) {\(e^{-i\theta}\)};
    \draw[black!60] (0.55,0) arc (0:50:0.55);
    \node[right,black!70] at (0.5,0.22) {\(\theta\)};
    \node at (0,-2.5) {复平面：两个特征值就位};
  \end{scope}
\end{tikzpicture}

\emph{左边转 \(\theta=50^\circ\)，每个向量都离开了自己张成的直线，实特征向量无处安放。右边把同一个变换记在复标量上：\(R_\theta\) 的特征值是单位圆上的一对共轭点 \(e^{\pm i\theta}\)，幅角正是转过的角度。\(\theta=90^\circ\) 时它们就是 \(i\) 与 \(-i\)。}
\end{center}

一般的平面旋转

\[
R_\theta=\begin{bmatrix}\cos\theta&-\sin\theta\\\sin\theta&\cos\theta\end{bmatrix}
\]

特征多项式为 \(\lambda^2-2\lambda\cos\theta+1\)，根是 \(e^{\pm i\theta}\)，特征向量可取 \((1,\mp i)^{\trans}\)。验证 \(\theta=90^\circ\) 的情形：\(A(1,-i)^{\trans}=(i,1)^{\trans}=i\,(1,-i)^{\trans}\)。两维实平面上必须用矩阵描述的旋转，在复方向上缩成了一次标量乘法。

\subsection{代数封闭是要付的最小代价}

有人会觉得这只是换个记号躲开麻烦。实际上扩到 \(\C\) 买到的东西相当具体：代数基本定理保证任何 \(n\) 次多项式在复数中恰有 \(n\) 个根（按重数计），所以任何 \(n\times n\) 复方阵的特征多项式一定分裂，特征值一定存在。在 \(\R\) 上这个保证不成立，上面的旋转就是反例。

而且这笔账付一次就够。扩到 \(\C\) 之后不会再冒出``复系数多项式无解''的新缺口——\(\C\) 已经代数封闭，不需要继续造更大的数域。这就是把复数称为最小代价的理由：它恰好补齐了缺的那一块，没有多要。

实矩阵还额外附赠一条对称性。特征多项式的系数都是实数，取共轭不改变方程，所以非实特征值必成共轭对出现；相应地，若 \(Av=\lambda v\)，则 \(A\overline v=\overline\lambda\,\overline v\)。一般复矩阵没有这个保证，\(\diag(i,2)\) 的谱就不成对。

模长与幅角的分工在迭代时看得最清楚。若 \(\lambda=re^{i\theta}\)，则 \(\lambda^k=r^ke^{ik\theta}\)：每走一步，相位推进 \(\theta\)，半径乘上 \(r\)。所以复特征值读作``一个振荡模态'' ——幅角定振荡频率，模长定它是衰减、维持还是发散。

\begin{center}
\begin{tikzpicture}[x=1cm,y=1cm,font=\footnotesize,>=stealth]
  \draw[black!20] (0,0) circle (1.7);
  \draw[black!35] (-1.95,0) -- (2.05,0);
  \draw[black!35] (0,-1.95) -- (0,2.05);
  \foreach \k in {0,...,15} {
    \draw[blue!60!black,thick] ({1.7*pow(0.86,\k)*cos(38*\k)},{1.7*pow(0.86,\k)*sin(38*\k)})
      -- ({1.7*pow(0.86,\k+1)*cos(38*(\k+1))},{1.7*pow(0.86,\k+1)*sin(38*(\k+1))});
    \fill[blue!60!black] ({1.7*pow(0.86,\k)*cos(38*\k)},{1.7*pow(0.86,\k)*sin(38*\k)}) circle (1.1pt);
  }
  \node at (0,-2.5) {\(r=0.86<1\)：旋进原点};
  \begin{scope}[shift={(6.0,0)}]
    \draw[black!20] (0,0) circle (0.75);
    \draw[black!35] (-1.95,0) -- (2.05,0);
    \draw[black!35] (0,-1.95) -- (0,2.05);
    \foreach \k in {0,...,11} {
      \draw[red!60!black,thick] ({0.75*pow(1.12,\k)*cos(38*\k)},{0.75*pow(1.12,\k)*sin(38*\k)})
        -- ({0.75*pow(1.12,\k+1)*cos(38*(\k+1))},{0.75*pow(1.12,\k+1)*sin(38*(\k+1))});
      \fill[red!60!black] ({0.75*pow(1.12,\k)*cos(38*\k)},{0.75*pow(1.12,\k)*sin(38*\k)}) circle (1.1pt);
    }
    \node at (0,-2.5) {\(r=1.12>1\)：旋出去};
  \end{scope}
\end{tikzpicture}

\emph{同一个幅角 \(\theta=38^\circ\)，两种模长。灰圈是单位圆。每一步转过的角度相同，所以两条轨迹的``转速''一致；决定去留的只有半径的乘子。线性系统的稳定性判据由此而来：离散时间看特征值是否全在单位圆内，连续时间看实部是否全为负。}
\end{center}

\subsection{转置必须换成共轭转置}

把 \(\R^n\) 换成 \(\C^n\) 后，加法、矩阵乘法、消元、秩、行列式几乎原样保留。唯一不能照抄的是内积，而这一处不改就寸步难行。

天真地沿用 \(x^{\trans}x\)，取 \(x=(1,i)^{\trans}\)：

\[
x^{\trans}x=1^2+i^2=0.
\]

一个非零向量长度为零。正定性一垮，正交、投影、最小二乘、能量守恒全部跟着塌。修法只有一处改动，在第一个变量上取共轭：

\[
\ip{x}{y}=x^{\conj}y=\sum_j\overline{x_j}\,y_j,
\qquad
x^{\conj}=\overline{x}^{\,\trans}.
\]

于是 \(\norm{x}^2=x^{\conj}x=\sum_j\abs{x_j}^2\)，非负，且只在 \(x=0\) 时取零。代价是内积不再对称，而是共轭对称：\(\ip{y}{x}=\overline{\ip{x}{y}}\)，并且它对第一个变量共轭线性、对第二个变量线性。有的教材把共轭放在第二个变量上，结论等价，但同一份推导里不能中途换边。

共轭转置 \(A^{\conj}\) 的作用可以脱离坐标来说。它是 \(A\) 关于这个内积的伴随：

\[
\ip{Ax}{y}=(Ax)^{\conj}y=x^{\conj}A^{\conj}y=\ip{x}{A^{\conj}y}.
\]

把变换从内积的一侧搬到另一侧，搬过去以后变成哪个矩阵，答案就是 \(A^{\conj}\)。实矩阵里共轭什么也没做，\(A^{\conj}\) 退回 \(A^{\trans}\)，所以以前那套写法一直没出错——它只是复情形的一个特例。

\subsection{酉矩阵接管正交矩阵的位置}

有了正确的内积，几何几乎全额恢复。\(x\perp y\) 记作 \(x^{\conj}y=0\)；单位向量 \(q\) 上的投影是 \(p=q(q^{\conj}x)\)；列标准正交的矩阵满足 \(Q^{\conj}Q=I\)。方阵情形单独给个名字：若

\[
U^{\conj}U=UU^{\conj}=I,
\]

称 \(U\) 为酉矩阵（unitary matrix）。它保内积，因为 \(\ip{Ux}{Uy}=x^{\conj}U^{\conj}Uy=x^{\conj}y\)，于是长度、夹角、总能量一并保住，而且求逆不用解方程，\(U^{-1}=U^{\conj}\)。

这两点合起来说明酉换基在数值上是安全的：矩阵的条件数为 \(1\)，误差进去多大出来还是多大，不会被放大；变回原坐标只需一次共轭转置。后面 Fourier 变换之所以能当作日常工具反复用，靠的正是它是一次酉换基。举个小例子，

\[
U=\frac1{\sqrt2}\begin{bmatrix}1&i\\i&1\end{bmatrix}
\]

的两列范数为一且相互正交，是酉矩阵。它把两个坐标的幅值与相位搅在一起，却不动向量的二范数。

\subsection{哪些直觉不能顺手带过来}

复数没有凭空造出新自由度。一个 \(\C^n\) 中的向量拆成实部与虚部就是 \(\R^{2n}\) 中的向量，一个复线性算子也有对应的实分块矩阵。选复坐标的理由不是能装更多信息，而是旋转、相位和共轭对称在这套坐标里直接可见，写出来的式子短。

但有两处要留神。第一，共轭映射 \(z\mapsto\overline z\) 只是实线性的：\(\overline{iz}=-i\,\overline z\)，它不能写成复矩阵乘法。凡是含共轭、取模、取实部的操作，都得先确认自己在哪个数域上讨论线性。第二，也是更常被误读的一处：特征多项式分裂不等于矩阵可对角化。\(\begin{bmatrix}1&1\\0&1\end{bmatrix}\) 在 \(\C\) 上的特征值仍是重根 \(1\)，独立特征向量仍然只有一个。复数解决的是``根存在吗''，解决不了``特征向量够不够'' ——\hyperref[sec:03-Linear-Algebra/05-Eigenvalues-and-Eigenvectors/03]{``对角化失败与 Jordan 形''}讨论的那些麻烦一个也没消失。

那么，在复空间里哪一类矩阵能保证有一组标准正交的特征基？实数情形的答案是对称矩阵，条件是 \(A^{\trans}=A\)。现在转置已经换成了共轭转置，下一节先把这个条件照样改写，再看看改写之后的谱定理还剩下什么、多出什么。
